\uuid{ticc}
\exo7id{194}
\titre{exo7 194}
\auteur{ridde}
\organisation{exo7}
\datecreate{1999-11-01}
\isIndication{false}
\isCorrection{false}
\chapitre{Injection, surjection, bijection}
\sousChapitre{Injection, surjection}
\cpge{mpsi-mp2i.raisonnement-vocabulaire-ensembliste.b, pcsi.raisonnement-vocabulaire-ensembliste.b, ptsi.raisonnement-vocabulaire-ensembliste.b, mpsi-mp2i.raisonnement-vocabulaire-ensembliste.c, pcsi.raisonnement-vocabulaire-ensembliste.d, ptsi.raisonnement-vocabulaire-ensembliste.d}
\module{Algèbre}
\niveau{L1}
\difficulte{}

\contenu{
\texte{
Soit $f : X \rightarrow Y$. Montrer que
}
\begin{enumerate}
    \item \question{$\forall B \subset Y \, \, f (f^{-1} (B)) = B \cap f (X)$.}
    \item \question{$f$ est surjective ssi $\forall B \subset Y \, \, f (f^{-1} (B)) = B $.}
    \item \question{$f$ est injective ssi $\forall A \subset X \, \, f^{-1} (f(A)) = A $.}
    \item \question{$f$ est bijective ssi $\forall A \subset X \, \,f (\complement A) = \complement f (A).$}
\end{enumerate}
}
